\section{Ablation Study}

This section analyzes the contributions of biomedical semantics, hypergraph structure, and medical knowledge constraints through controlled architecture comparisons rather than single-feature-group removal. As noted in Section~5, leave-one-feature-group ablations are retained in the appendix as complementary feature-coverage diagnostics, but are not used as the primary ablation evidence because high multicollinearity among knowledge features leads to automatic compensation.

\begin{table}[t]
\centering
\small
\caption{Structural ablation: component contributions via controlled architecture comparisons on held-out BioASQ. Pairwise bootstrap significance against the second method in each pair.}
\label{tab:structural-ablation}
\begin{tabular}{lrrrrrr}
\toprule
Comparison & Rec@10 & MRR@10 & nDCG@10 & $\Delta$Rec@10 & $\Delta$MRR@10 & $\Delta$nDCG@10 \\
\midrule
(1) Retrieval LTR & 0.5208 & 0.7791 & 0.6282 & -- & -- & -- \\
(1) Semantic (no graph) & 0.5309 & 0.7872 & 0.6436 & +0.0101** & +0.0080* & +0.0153** \\
\midrule
(2) Semantic (no graph) & 0.5309 & 0.7872 & 0.6436 & -- & -- & -- \\
(2) Hypergraph (no med) & 0.5266 & 0.7782 & 0.6362 & $-$0.0044 & $-$0.0089** & $-$0.0074** \\
\midrule
(3) Hypergraph (no med) & 0.5266 & 0.7782 & 0.6362 & -- & -- & -- \\
(3) KCH-MedRank & 0.5329 & 0.7867 & 0.6433 & +0.0063** & +0.0085** & +0.0071** \\
\midrule
(4) Pairwise graph & 0.5306 & 0.7866 & 0.6409 & -- & -- & -- \\
(4) KCH-MedRank & 0.5329 & 0.7867 & 0.6433 & +0.0023 & +0.0001 & +0.0024 \\
\bottomrule
\end{tabular}
\end{table}


\subsection{Biomedical Semantics Provide a Strong Independent Gain}

Comparison (1) in Table~\ref{tab:structural-ablation} isolates the contribution of biomedical semantic features above retrieval signals alone. Adding the precomputed biomedical semantic score and its rank-based feature improves Recall@10 from 0.5208 to 0.5309 ($\Delta = +0.0101$, $p = 0.0006$) and nDCG@10 from 0.6282 to 0.6436 ($\Delta = +0.0153$, $p = 0.0002$). MRR@10 improves from 0.7791 to 0.7872 ($\Delta = +0.0080$, $p = 0.029$). These results confirm that specialized biomedical dense retrieval provides evidence ranking signals that are not fully captured by general-purpose lexical and dense retrieval alone.

\subsection{The Hypergraph Alone Introduces Noise}

Comparison (2) isolates the effect of the local hypergraph structure without medical knowledge constraints. This variant builds the full hypergraph topology (candidate membership edges, rank-band edges, entity-type concept edges, and shared-entity passage edges) but replaces all biomedical MeSH, entity, and PrimeKG features with zeros. The resulting model underperforms the semantic-only LambdaMART baseline: nDCG@10 drops by 0.0074 ($p = 0.0002$) and MRR@10 drops by 0.0089 ($p = 0.003$). Recall@10 drops by 0.0044 (ns, trending negative). This indicates that the raw hypergraph topology, without medically grounded edge weights and concept annotations, introduces structural noise that degrades ranking quality.

\subsection{Medical Knowledge Constraints Transform the Hypergraph from Noise to Signal}

Comparison (3) is the central structural ablation. It compares KCH-MedRank against the same hypergraph architecture with all medical knowledge features removed. Adding biomedical entity features, exact MeSH features, MeSH hierarchy features, and PrimeKG relation features improves Recall@10 by 0.0063 ($p = 0.0028$), MRR@10 by 0.0085 ($p = 0.0042$), and nDCG@10 by 0.0071 ($p = 0.0002$). All three top-10 metrics are statistically significant.

\begin{table}[t]
\centering
\small
\caption{Interpretation of structural ablation results. Knowledge$\times$Structure interaction is the difference between the hypergraph with and without medical knowledge constraints.}
\label{tab:structural-ablation-interpretation}
\begin{tabular}{lrrrc}
\toprule
Component & $\Delta$Rec@10 & $\Delta$MRR@10 & $\Delta$nDCG@10 & Role \\
\midrule
Biomedical semantics & +0.0101** & +0.0080* & +0.0153** & independent gain \\
Hypergraph alone (no med) & $-$0.0044 & $-$0.0089** & $-$0.0074** & introduces noise \\
Medical knowledge constraint & +0.0063** & +0.0085** & +0.0071** & transforms noise to signal \\
Knowledge $\times$ Structure & +0.0107 & +0.0174 & +0.0145 & net interaction effect \\
\bottomrule
\end{tabular}
\end{table}


Table~\ref{tab:structural-ablation-interpretation} summarizes the interpretation. The hypergraph structure alone degrades ranking, but medical knowledge constraints reverse the effect, yielding a net positive interaction. The Knowledge$\times$Structure interaction effect (the difference between the hypergraph with and without medical knowledge) is $+0.0107$ for Recall@10 and $+0.0174$ for MRR@10. This provides the primary ablation evidence that medical knowledge constraints are necessary for the hypergraph to contribute positively to evidence ranking.

\subsection{Hypergraph Structure Marginally Outperforms Pairwise Graph}

Comparison (4) evaluates whether the n-ary hyperedge representation provides benefits over a pairwise graph decomposition that breaks each multi-node hyperedge into anchor-based pairwise edges. KCH-MedRank improves Recall@10 over pairwise graph LTR by 0.0023 and nDCG@10 by 0.0024, but these differences are not statistically significant at $k=10$ ($p=0.205$, paired bootstrap). At $k=5$, Recall@5 improves by 0.0054 ($p=0.005$, paired bootstrap) and nDCG@5 by 0.0042, confirming a modest but statistically significant early-rank benefit.

\paragraph{Stratified Analysis.} To understand when the hypergraph structure matters most, we stratify the test set by the structural characteristics of each question. Table~\ref{tab:hypergraph-vs-pairwise-stratified} presents per-stratum paired bootstrap comparisons at $k=5$ and $k=10$.

\begin{table}[t]
\centering
\small
\caption{Stratified paired bootstrap analysis of KCH-MedRank vs. Pairwise Graph LTR, grouped by gold passage structural characteristics at $k=5$ and $k=10$. ``MeSH Overlap'' refers to whether the gold passage shares any MeSH terms with the question. ``Entity Overlap'' refers to shared biomedical entities. $p$-values are two-sided from 10,000 paired bootstrap iterations.}
\label{tab:hypergraph-vs-pairwise-stratified}
\begin{tabular}{lrrrrrl}
\toprule
Subset & $n$ & KCH Rec@5 & PW Rec@5 & $\Delta$ & 95\% CI & $p$ \\
\midrule
\multicolumn{7}{c}{\textbf{Recall@5}} \\
\midrule
Overall & 943 & 0.4210 & 0.4156 & +0.0054 & [+0.0016, +0.0099] & 0.005** \\
\midrule
MeSH Overlap = 0 & 366 & 0.4114 & 0.4016 & +0.0098 & [+0.0021, +0.0196] & 0.006** \\
MeSH Overlap $\geq$ 1 & 577 & 0.4271 & 0.4245 & +0.0026 & [--0.0010, +0.0065] & 0.159 \\
Entity Overlap = 0 & 649 & 0.4177 & 0.4132 & +0.0046 & [--0.0002, +0.0102] & 0.062 \\
Entity Overlap $\geq$ 1 & 294 & 0.4283 & 0.4211 & +0.0072 & [+0.0014, +0.0136] & 0.015* \\
1 gold passage & 204 & 0.5637 & 0.5539 & +0.0098 & [+0.0000, +0.0245] & 0.277 \\
$\geq$ 2 gold passages & 739 & 0.3816 & 0.3775 & +0.0042 & [+0.0006, +0.0079] & 0.023* \\
\midrule
\multicolumn{7}{c}{\textbf{Recall@10}} \\
\midrule
Overall & 943 & 0.5329 & 0.5306 & +0.0023 & [--0.0011, +0.0064] & 0.205 \\
\midrule
MeSH Overlap = 0 & 366 & 0.4754 & 0.4722 & +0.0032 & [--0.0021, +0.0102] & 0.306 \\
MeSH Overlap $\geq$ 1 & 577 & 0.5693 & 0.5676 & +0.0017 & [--0.0024, +0.0067] & 0.471 \\
Entity Overlap = 0 & 649 & 0.5205 & 0.5162 & +0.0043 & [--0.0002, +0.0096] & 0.066 \\
Entity Overlap $\geq$ 1 & 294 & 0.5603 & 0.5623 & --0.0020 & [--0.0068, +0.0025] & 0.378 \\
1 gold passage & 204 & 0.5833 & 0.5735 & +0.0098 & [+0.0000, +0.0245] & 0.273 \\
$\geq$ 2 gold passages & 739 & 0.5189 & 0.5187 & +0.0002 & [--0.0027, +0.0031] & 0.867 \\
\bottomrule
\end{tabular}
\end{table}


The stratification reveals two key patterns. First, the hypergraph advantage is largest when the gold passage has \textit{no direct MeSH overlap} with the question ($\Delta_{\text{Recall@5}} = +0.0098$, $p=0.006$). This makes structural sense: when direct MeSH overlap exists, the pairwise graph decomposition already captures most of the alignment signal through shared-concept edges; when absent, the hypergraph's n-ary structure enables indirect paths through shared-entity clusters and MeSH hierarchy ancestors that the pairwise decomposition fragments. Second, questions with $\geq 2$ gold passages also benefit significantly from the hypergraph ($\Delta_{\text{Recall@5}} = +0.0042$, $p=0.023$), consistent with multi-evidence scenarios where the hypergraph's shared-entity hyperedges propagate ranking signal across multiple gold passages.

At $k=10$, no subset reaches significance, but the effect direction is consistent ($\Delta$ positive in all but one stratum). The single-gold-passage subset shows the largest nominal $\Delta$ (+0.0098), though the smaller sample ($n=204$) limits power. This confirms that the hypergraph advantage is concentrated in the early-rank region and in cases where indirect knowledge signals compensate for missing direct concept alignment.

\subsection{Summary of Ablation Evidence}

Taken together, the structural ablations support three conclusions. First, biomedical semantic features are the strongest single component addition to retrieval-based learning-to-rank. Second, the hypergraph topology alone is harmful without medical concept annotations. Third, medical knowledge constraints are necessary and effective: they reverse the hypergraph's noise into a statistically significant signal, producing a net Knowledge$\times$Structure interaction that improves all top-10 evidence retrieval metrics. The leave-one-feature-group diagnostics in the appendix complement this analysis by showing that individual knowledge feature groups are highly redundant with each other, further supporting the interpretation that medical constraints work as an integrated block rather than as independent additive components.
